% ============================================================================
% GLOSARIO (Optativo)
% ============================================================================
% Definición de términos técnicos y abreviaturas utilizadas en el documento.
% Si no deseas incluir glosario, puedes dejar este archivo vacío o no incluirlo en main.tex.
% ============================================================================

\chapter*{GLOSARIO}
\addcontentsline{toc}{chapter}{GLOSARIO}

% Se recomienda usar un entorno de descripción (description list) para alinear términos
\begin{description}[leftmargin=3cm, style=nextline]
    \item[API] (\textit{Application Programming Interface}): Interfaz de programación de aplicaciones. Conjunto de definiciones y protocolos para el desarrollo e integración de software.
    
    \item[CSS] (\textit{Cascading Style Sheets}): Hojas de estilo en cascada, utilizadas para definir la presentación visual de documentos estructurados en HTML o XML.
    
    \item[HTML] (\textit{HyperText Markup Language}): Lenguaje de marcado de hipertexto. Es el estándar para la estructura y contenido de las páginas web.
    
    \item[LaTeX] Sistema de preparación de documentos para la composición tipográfica de alta calidad, ampliamente utilizado en textos científicos y académicos.
    
    \item[REST] (\textit{Representational State Transfer}): Transferencia de estado representacional. Estilo arquitectónico para el diseño de servicios web.
    
    \item[UML] (\textit{Unified Modeling Language}): Lenguaje de modelado unificado. Estándar visual para especificar, construir y documentar sistemas de software.
    
    \item[UTEM] Universidad Tecnológica Metropolitana.
\end{description}
